% ============================================================
%  Taller 1 — Filtros de Denoising en Imágenes MRI
%  Pontificia Universidad Javeriana
%  Procesamiento de Imágenes Médicas
%
%  Compilar con: pdflatex taller1_filtrado.tex
%  (dos pasadas para referencias cruzadas)
%
%  Estructura de carpetas esperada desde la raíz del proyecto:
%    output/comparison_results/   ← imágenes _wiener_comparison y _single_comparison
%    output/report/               ← imágenes _comparison (4 filtros) y summary_axial.png
% ============================================================

\documentclass[12pt, a4paper]{article}

% ── Codificación y lenguaje ────────────────────────────────
\usepackage[utf8]{inputenc}
\usepackage[T1]{fontenc}
\usepackage[spanish, es-tabla]{babel}

% ── Tipografía y geometría ─────────────────────────────────
\usepackage{lmodern}
\usepackage[margin=2.5cm]{geometry}
\usepackage{microtype}
\usepackage{parskip}
\setlength{\parindent}{0pt}
\setlength{\parskip}{6pt}

% ── Matemáticas ────────────────────────────────────────────
\usepackage{amsmath, amssymb, amsthm}

% ── Figuras y subfiguras ───────────────────────────────────
\usepackage{graphicx}
\usepackage{subcaption}
\usepackage{float}
\usepackage[export]{adjustbox}

% ── Tablas ────────────────────────────────────────────────
\usepackage{booktabs}
\usepackage{tabularx}
\usepackage{multirow}
\usepackage{xcolor}
\usepackage{colortbl}
\usepackage{array}

% ── Código fuente ─────────────────────────────────────────
\usepackage{listings}
\lstset{
  basicstyle=\ttfamily\small,
  keywordstyle=\color{blue!80!black}\bfseries,
  commentstyle=\color{green!50!black},
  stringstyle=\color{red!70!black},
  numbers=left,
  numberstyle=\tiny\color{gray},
  stepnumber=1,
  frame=single,
  breaklines=true,
  showstringspaces=false,
  language=Python,
  backgroundcolor=\color{gray!8},
  rulecolor=\color{gray!40},
}

% ── Hipervínculos ─────────────────────────────────────────
\usepackage[colorlinks=true, linkcolor=blue!70!black,
            citecolor=green!50!black, urlcolor=blue!70!black]{hyperref}

% ── Encabezados y pies de página ──────────────────────────
\usepackage{fancyhdr}
\pagestyle{fancy}
\fancyhf{}
\fancyhead[L]{\small Pontificia Universidad Javeriana $\mid$ Procesamiento de Imágenes Médicas}
\fancyhead[R]{\small Taller 1 — Filtros de Denoising}
\fancyfoot[C]{\thepage}
\renewcommand{\headrulewidth}{0.4pt}

% ── Cajas de información ───────────────────────────────────
\usepackage{tcolorbox}
\tcbuselibrary{skins, breakable}
\tcbset{
  infostyle/.style={
    colback=blue!5, colframe=blue!40!black,
    fonttitle=\bfseries, breakable, enhanced
  },
  warnstyle/.style={
    colback=orange!8, colframe=orange!70!black,
    fonttitle=\bfseries, breakable, enhanced
  },
  codestyle/.style={
    colback=gray!6, colframe=gray!50,
    fonttitle=\bfseries, breakable, enhanced
  }
}

% ── Títulos de sección personalizados ─────────────────────
\usepackage{titlesec}
\titleformat{\section}{\large\bfseries\color{blue!70!black}}{}{0em}
  {\thesection\quad}[\titlerule]
\titleformat{\subsection}{\normalsize\bfseries\color{blue!50!black}}
  {\thesubsection\quad}{0em}{}
\titleformat{\subsubsection}{\normalsize\itshape\bfseries}
  {\thesubsubsection\quad}{0em}{}

% ── Rutas base de imágenes ────────────────────────────────
% Ajustar si el .tex no está en la raíz del proyecto
\graphicspath{
  {output/report/}
  {output/comparison_results/}
}

% ── Metadatos del documento ───────────────────────────────
\title{
  \vspace{-1cm}
  {\large Pontificia Universidad Javeriana}\\[0.3em]
  {\large Procesamiento de Imágenes Médicas}\\[1em]
  \rule{\linewidth}{1.5pt}\\[0.5em]
  {\LARGE\bfseries Taller 1 — Filtros de Denoising\\[0.2em]
   en Imágenes de Resonancia Magnética}\\[0.3em]
  \rule{\linewidth}{1.5pt}
}
\author{
  Abel Albuez\quad$\mid$\quad Victoria Acero\quad$\mid$\quad Santiago Gil\\[0.4em]
  \small\texttt{aa-albuezs@javeriana.edu.co}
  $\;\cdot\;$
  \texttt{jv.acero@javeriana.edu.co}
  $\;\cdot\;$
  \texttt{santiago\_gil@javeriana.edu.co}
}
\date{Bogotá D.C., Colombia — Marzo 2026}

% ============================================================
\begin{document}
% ============================================================

\maketitle
\thispagestyle{fancy}

% ── Resumen ───────────────────────────────────────────────
\begin{tcolorbox}[infostyle, title={Resumen}]
Este documento reporta la implementación y análisis experimental de tres filtros clásicos
de denoising aplicados a imágenes de resonancia magnética (MRI): el \textbf{Filtro de
Mediana} (MF), el \textbf{Filtro de Mediana Adaptativo} (AMF) y el \textbf{Filtro de
Wiener Adaptativo} (AWF). Los tres filtros son métodos del dominio espacial implementados
en Python con soporte de la biblioteca \emph{Insight Toolkit} (ITK) para entrada/salida
de volúmenes NIfTI, y NumPy/SciPy para la lógica de filtrado. Los experimentos evalúan
el comportamiento de cada filtro ante distintos niveles de ruido Rician/Gaussiano mixto
sobre volúmenes cerebrales T1 del modelo BrainWeb (pn3, pn5 y pn9) con campo de
inhomogeneidad rf20.
\end{tcolorbox}

\vspace{0.5em}
\tableofcontents
\newpage

% ============================================================
\section{Introducción}
% ============================================================

El ruido es una de las principales fuentes de degradación en imágenes de resonancia
magnética. Durante la adquisición, el hardware del escáner introduce fluctuaciones
térmicas en la señal, que se manifiestan como ruido Rician en magnitud —asintóticamente
equivalente a ruido Gaussiano para relaciones señal-ruido moderadas. Adicionalmente,
errores de transmisión y defectos de sensor pueden generar ruido de impulso (sal y
pimienta). La presencia de ruido afecta directamente la calidad diagnóstica y compromete
el desempeño de etapas posteriores como segmentación, registro y extracción de
características.

El objetivo de este taller es implementar y comparar tres filtros de denoising del dominio
espacial, siguiendo la taxonomía propuesta por \cite{ali2018}:

\begin{table}[H]
\centering
\renewcommand{\arraystretch}{1.3}
\begin{tabularx}{\textwidth}{lllX}
\toprule
\textbf{Sección} & \textbf{Filtro} & \textbf{Dominio} & \textbf{Estado} \\
\midrule
2.2.1 & Filtro de Mediana (MF) & Espacial, no lineal & Implementado — resultados completos \\
2.2.2 & Filtro de Mediana Adaptativo (AMF) & Espacial, no lineal & Implementado — resultados completos \\
2.1   & Filtro de Wiener Adaptativo (AWF)  & Espacial, estadístico & Implementado — resultados completos \\
\bottomrule
\end{tabularx}
\caption{Filtros implementados en este taller (Ali 2018).}
\end{table}

Los tres filtros se aplican sobre volúmenes BrainWeb \cite{brainweb} con tres niveles de
ruido (pn3, pn5, pn9) y se evalúan cualitativamente mediante comparaciones visuales en
los tres planos anatómicos (axial, sagital, coronal).

% ============================================================
\section{Marco Teórico}
% ============================================================

\subsection{Tipos de ruido en imágenes MRI}

Las imágenes de resonancia magnética son susceptibles principalmente a dos tipos de
ruido con propiedades estadísticas diferenciadas:

\begin{table}[H]
\centering
\renewcommand{\arraystretch}{1.4}
\begin{tabularx}{\textwidth}{lXX}
\toprule
\textbf{Tipo} & \textbf{Característica} & \textbf{Distribución matemática} \\
\midrule
Gaussiano / Rician &
Afecta todos los vóxeles con pequeñas variaciones continuas. Originado en el ruido
térmico del hardware del escáner. &
$p(z) = \dfrac{1}{\sigma\sqrt{2\pi}}\, e^{-\frac{(z-\mu)^2}{2\sigma^2}}$ \\[1em]
Sal y pimienta (impulso) &
Vóxeles aislados con valores extremos (0 o MAX). Causado por errores de sensor o
transmisión. &
$P(z) = \begin{cases} p_a & z = a \\ p_b & z = b \\ 0 & \text{otro caso} \end{cases}$ \\[1em]
Mixto &
Combinación simultánea de Gaussiano e impulso. Escenario más exigente para los filtros. &
Superposición de ambas distribuciones. \\
\bottomrule
\end{tabularx}
\caption{Tipos de ruido considerados en los experimentos.}
\end{table}

\subsection{Filtro de Mediana (MF)}

El Filtro de Mediana es un filtro no lineal del dominio espacial. Su operación consiste en
reemplazar el valor de cada vóxel por la \textbf{mediana estadística} de los valores en
una vecindad rectangular $S_{xy}$ centrada en dicho vóxel:

\begin{equation}
  \hat{f}(x,y) = \operatorname{mediana}\bigl\{g(s,t)\bigr\},\quad (s,t) \in S_{xy}
  \label{eq:median}
\end{equation}

donde $g(x,y)$ es la imagen corrupta y $\hat{f}(x,y)$ la imagen restaurada. A diferencia
del filtro de media, la mediana es inherentemente robusta ante valores extremos aislados
(ruido de impulso) porque es una estadística de orden: los valores más alejados del centro
de la distribución no influyen en el resultado final.

\textbf{Parámetro clave}: el radio de la ventana (\texttt{radius}). Un radio mayor implica
mayor suavizado y mayor reducción de ruido a costa de potencial pérdida de bordes y
detalles finos. En la implementación 3D del presente trabajo se utiliza una vecindad cúbica
$(2r+1)^3$, siendo $r=1$ la configuración por defecto (ventana $3\times3\times3 = 27$
vóxeles).

\subsection{Filtro de Mediana Adaptativo (AMF)}

El Filtro de Mediana Adaptativo \cite{lin2013} extiende el filtro clásico mediante un
\textbf{mecanismo de ventana dinámica}. En lugar de usar una vecindad de tamaño fijo,
adapta el tamaño de la ventana vóxel a vóxel según la densidad local de ruido.

El algoritmo opera en dos niveles de decisión sobre cada píxel $(x,y)$, donde se definen:
$Z_{\min}$ = mínimo de la vecindad, $Z_{\max}$ = máximo, $Z_{\text{med}}$ = mediana,
$Z_{xy}$ = valor del vóxel actual, $S_{\max}$ = tamaño máximo de ventana permitido.

\begin{tcolorbox}[codestyle, title={Algoritmo AMF — Dos niveles de decisión}]
\textbf{Nivel A} — ¿Es la mediana confiable (no es ruido)?
\begin{enumerate}
  \item Calcular $L_{11} = Z_{\text{med}} - Z_{\min}$ y $L_{12} = Z_{\text{med}} - Z_{\max}$
  \item Si $L_{11} > 0$ \textbf{AND} $L_{12} < 0$ $\Rightarrow$ mediana confiable, ir al \textbf{Nivel B}
  \item Si no $\Rightarrow$ ampliar ventana ($3 \to 5 \to 7 \to \ldots$) y repetir Nivel A
  \item Si ventana $> S_{\max}$ $\Rightarrow$ output $= Z_{\text{med}}$
\end{enumerate}
\textbf{Nivel B} — ¿Es el vóxel central ruido o señal limpia?
\begin{enumerate}
  \item Calcular $L_{21} = Z_{xy} - Z_{\min}$ y $L_{22} = Z_{xy} - Z_{\max}$
  \item Si $L_{21} > 0$ \textbf{AND} $L_{22} < 0$ $\Rightarrow$ vóxel limpio $\Rightarrow$ output $= Z_{xy}$ (conservar)
  \item Si no $\Rightarrow$ vóxel es ruido $\Rightarrow$ output $= Z_{\text{med}}$ (reemplazar)
\end{enumerate}
\end{tcolorbox}

La clave conceptual del algoritmo es que un vóxel solo se modifica si hay \textit{evidencia
matemática} de que es un valor extremo no representativo del tejido local. Esto le permite
preservar bordes y detalles estructurales que el filtro clásico de ventana fija tiende a
suavizar.

\subsection{Filtro de Wiener Adaptativo (AWF)}

El Filtro de Wiener Adaptativo es un estimador de mínimo error cuadrático medio (MMSE)
que adapta su comportamiento según las \textbf{estadísticas locales de media y varianza}
calculadas en una ventana rectangular $M \times N$ alrededor de cada vóxel
\cite{ali2018}. Opera en el dominio espacial (no en el dominio de la frecuencia como el
filtro de Wiener clásico).

\subsubsection{Formulación matemática}

La fórmula de restauración por vóxel es:

\begin{equation}
  \hat{f}(x,y) = \mu + \frac{\max\bigl(0,\;\sigma^2_{\text{local}} - \nu^2\bigr)}
                             {\max\bigl(\sigma^2_{\text{local}},\;\varepsilon\bigr)}
                 \cdot \bigl[g(x,y) - \mu\bigr]
  \label{eq:wiener}
\end{equation}

donde:
\begin{itemize}
  \item $g(x,y)$ — vóxel de la imagen ruidosa de entrada
  \item $\mu$ — media local en la ventana $M \times N$ centrada en $(x,y)$
  \item $\sigma^2_{\text{local}}$ — varianza local en la misma ventana: $\sigma^2_{\text{local}} = E[X^2] - E[X]^2$
  \item $\nu^2$ — varianza global del ruido (estimada automáticamente como promedio de todas las varianzas locales)
  \item $\varepsilon$ — épsilon de máquina ($\approx 2.2 \times 10^{-16}$), evita división por cero
\end{itemize}

\subsubsection{Comportamiento adaptativo}

La fracción en la ecuación \eqref{eq:wiener} actúa como una \textbf{ganancia de Wiener}
que varía continua\-mente entre 0 y 1:

\begin{table}[H]
\centering
\renewcommand{\arraystretch}{1.4}
\begin{tabularx}{\textwidth}{lXX}
\toprule
\textbf{Condición local} & \textbf{Valor de la ganancia} & \textbf{Efecto sobre el vóxel} \\
\midrule
$\sigma^2_{\text{local}} \leq \nu^2$ (región plana/ruidosa) &
ganancia $= 0$ &
$\hat{f} = \mu$ — suavizado máximo \\
$\sigma^2_{\text{local}} \gg \nu^2$ (borde / estructura real) &
ganancia $\to 1$ &
$\hat{f} \approx g$ — señal original preservada \\
Caso intermedio &
ganancia $\in (0,1)$ &
mezcla ponderada entre $\mu$ y $g$ \\
\bottomrule
\end{tabularx}
\caption{Comportamiento adaptativo del AWF según la varianza local.}
\end{table}

\subsubsection{Estimación de la varianza de ruido $\nu^2$}

En ausencia de una medición independiente del ruido, $\nu^2$ se estima automáticamente
como el promedio de todas las varianzas locales del volumen:

\begin{equation}
  \nu^2 \approx \frac{1}{N_{\text{vóxeles}}} \sum_{x,y,z} \sigma^2_{\text{local}}(x,y,z)
  \label{eq:nu2}
\end{equation}

Esta estimación es razonable porque, en presencia de ruido, las varianzas locales están
infladas precisamente en $\nu^2$ respecto de las varianzas del tejido limpio; promediar
sobre todo el volumen proporciona un estimador estable.

% ============================================================
\section{Implementación}
% ============================================================

Los tres filtros se implementaron en Python 3 con la siguiente arquitectura común:

\begin{tcolorbox}[infostyle, title={Pipeline general (los tres filtros)}]
\textbf{Entrada:} Volumen NIfTI \texttt{.nii.gz} (BrainWeb) \\[0.3em]
\textbf{Paso 1 — Lectura ITK:} \texttt{itk.imread()} carga el volumen y lo convierte a
\texttt{itk.Image[itk.F, 3]} (float32, 3D).\\[0.3em]
\textbf{Paso 2 — Conversión:} \texttt{itk.GetArrayFromImage()} exporta un array NumPy
de forma $(Z, Y, X)$ float32.\\[0.3em]
\textbf{Paso 3 — Filtrado:} se aplica el algoritmo con NumPy/SciPy.\\[0.3em]
\textbf{Paso 4 — Reconstrucción ITK:} \texttt{itk.image\_from\_array()} + copia de
\texttt{Spacing/Origin/Direction} del volumen original.\\[0.3em]
\textbf{Salida:} Volumen NIfTI filtrado + PNG de comparación.
\end{tcolorbox}

\begin{tcolorbox}[warnstyle, title={Nota técnica — Rutas con caracteres no ASCII}]
El proyecto reside en una carpeta denominada \textit{imágenes médicas}. La capa C++ de
ITK no puede procesar rutas con caracteres no ASCII en Windows. Todos los scripts
resuelven esta limitación copiando el archivo a un directorio temporal con ruta ASCII,
procesándolo allí y trasladando el resultado a la ubicación final.
\end{tcolorbox}

\subsection{Filtro de Mediana — \texttt{src/median.py}}

La implementación usa \texttt{numpy.lib.stride\_tricks.sliding\_window\_view} para
vectorizar la extracción de vecindades sin un bucle Python explícito sobre los $\sim$7
millones de vóxeles del volumen T1:

\begin{enumerate}
  \item Extraer la vecindad $(2r+1)^3$ alrededor de cada vóxel mediante \texttt{sliding\_window\_view}.
  \item Eliminar el vóxel central (índice 13 en la vecindad aplanada de 27 elementos).
  \item Reemplazar el vóxel con \texttt{np.median()} de los 26 vecinos restantes.
\end{enumerate}

ITK se usa exclusivamente para lectura (\texttt{.nii.gz}) y escritura (\texttt{.nii}) del
volumen. Parámetro principal: \texttt{radius=1} ($3\times3\times3$, por defecto en
\texttt{run\_all.py}).

\subsection{Filtro de Mediana Adaptativo — \texttt{src/adaptive-median.py}}

El script ofrece dos variantes seleccionables con el flag \texttt{--no-itk}:

\subsubsection{Variante ITK (por defecto)}

Aproximación al algoritmo adaptativo usando primitivas nativas de ITK:
\begin{enumerate}
  \item \texttt{itk.MedianImageFilter} con radio $(1,1,0)$ — ventana $3\times3$ en el
    plano axial.
  \item \texttt{itk.GradientMagnitudeImageFilter} sobre la imagen mediana filtrada.
  \item Umbral al percentil 82 de la magnitud de gradiente: donde se supera el umbral
    (zona de borde), se restaura el vóxel original; en regiones homogéneas prevalece la
    mediana filtrada.
\end{enumerate}
Esta variante es significativamente más rápida que la implementación NumPy.

\subsubsection{Variante NumPy (algoritmo exacto de Ali 2018)}

Implementación literal del algoritmo de dos niveles descrito en la sección anterior.
Procesa el volumen corte axial por corte axial. Usa \texttt{np.pad(mode='reflect')} para
manejo de bordes. La ventana crece de 3 en dos hasta $S_{\max}$. La desventaja es el
costo computacional: el doble bucle Python por píxel hace esta variante impráctica para
uso rutinario en volúmenes 3D grandes.

\begin{table}[H]
\centering
\renewcommand{\arraystretch}{1.3}
\begin{tabular}{llp{6.5cm}}
\toprule
\textbf{Argumento} & \textbf{Valor por defecto} & \textbf{Descripción} \\
\midrule
\texttt{--max-window} & 7 & $S_{\max}$: tamaño máximo de ventana (impar) \\
\texttt{--no-itk} & False & Activa la implementación NumPy exacta \\
\texttt{--noise-type} & none & Tipo de ruido a inyectar: \texttt{none | salt\_pepper | gaussian | mixed} \\
\texttt{--noise-density} & 0.1 & Fracción de píxeles afectados (sal y pimienta) \\
\texttt{--noise-sigma} & 10.0 & Desviación estándar del ruido Gaussiano \\
\bottomrule
\end{tabular}
\caption{Parámetros del script \texttt{adaptive-median.py}.}
\end{table}

\subsection{Filtro de Wiener Adaptativo — \texttt{src/wiener.py}}

Implementación completamente 3D y vectorizada con \texttt{scipy.ndimage.uniform\_filter}:

\begin{lstlisting}[caption={Núcleo del Filtro de Wiener Adaptativo (src/wiener.py)}]
from scipy.ndimage import uniform_filter
import numpy as np

def adaptive_wiener_filter_3d(arr, window_size=5, noise_variance=None):
    arr_f64 = arr.astype(np.float64)

    # Paso 1: media local (filtro de caja deslizante)
    local_mean = uniform_filter(arr_f64, size=window_size, mode='reflect')

    # Paso 2: varianza local = E[X^2] - E[X]^2
    local_mean_sq  = uniform_filter(arr_f64**2, size=window_size, mode='reflect')
    local_variance = local_mean_sq - local_mean**2

    # Paso 3: estimacion automatica de la varianza de ruido
    if noise_variance is None:
        noise_variance = float(np.mean(local_variance))

    # Paso 4: ganancia de Wiener (clampeada a [0, 1])
    eps         = np.finfo(np.float64).eps
    numerator   = np.maximum(local_variance - noise_variance, 0.0)
    denominator = np.maximum(local_variance, eps)
    gain        = numerator / denominator

    # Paso 5: restauracion voxel a voxel
    filtered = local_mean + gain * (arr_f64 - local_mean)
    return filtered.astype(np.float32)
\end{lstlisting}

\texttt{uniform\_filter(size=k)} calcula la media sobre una ventana cúbica $k\times k\times k$
deslizante en las tres dimensiones simultáneamente en un único paso a nivel de C, lo que
evita completamente los bucles Python. El parámetro \texttt{mode='reflect'} refleja el
array en los bordes del volumen, eliminando el sesgo de los artefactos de borde.

\subsection{Ejecutor por lotes — \texttt{src/run\_all.py}}

El script \texttt{run\_all.py} itera sobre las tres imágenes BrainWeb (pn3, pn5, pn9),
ejecuta los tres filtros sobre cada una y genera los artefactos de salida descritos en la
Tabla~\ref{tab:outputs}.

\begin{table}[H]
\centering
\renewcommand{\arraystretch}{1.3}
\begin{tabularx}{\textwidth}{lX}
\toprule
\textbf{Directorio de salida} & \textbf{Contenido} \\
\midrule
\texttt{output/median/} & 3 volúmenes NIfTI filtrados (pn3, pn5, pn9) \\
\texttt{output/adaptive-median/} & 3 volúmenes NIfTI filtrados \\
\texttt{output/wiener/} & 3 volúmenes NIfTI filtrados \\
\texttt{output/comparison\_results/} & PNGs individuales de cada filtro por nivel de ruido \\
\texttt{output/report/} & PNGs de comparación completa (4 filtros) y resumen axial \\
\bottomrule
\end{tabularx}
\caption{Artefactos generados por \texttt{run\_all.py}.}
\label{tab:outputs}
\end{table}

Cada PNG de comparación completa muestra el siguiente layout:
\[
  \underbrace{\text{Original}}_{\text{col. 1}} \quad
  \underbrace{\text{Mediana (r=1)}}_{\text{col. 2}} \quad
  \underbrace{\text{Med. Adaptativo}}_{\text{col. 3}} \quad
  \underbrace{\text{Wiener Adaptativo}}_{\text{col. 4}}
\]
en tres filas correspondientes a los planos axial, sagital y coronal.

% ============================================================
\section{Dataset — Imágenes BrainWeb}
% ============================================================

Los experimentos utilizan el simulador de MRI BrainWeb \cite{brainweb}, un modelo
anatómico digital del cerebro humano adulto que permite generar imágenes T1 con
parámetros de ruido y artefacto de inhomogeneidad controlados.

\begin{table}[H]
\centering
\renewcommand{\arraystretch}{1.3}
\begin{tabular}{llll}
\toprule
\textbf{Imagen} & \textbf{Tipo} & \textbf{Ruido (\%)} & \textbf{Inhomogeneidad RF (\%)} \\
\midrule
\texttt{t1\_icbm\_normal\_1mm\_pn3\_rf20} & T1 ICBM & 3\% (bajo)  & 20\% \\
\texttt{t1\_icbm\_normal\_1mm\_pn5\_rf20} & T1 ICBM & 5\% (medio) & 20\% \\
\texttt{t1\_icbm\_normal\_1mm\_pn9\_rf20} & T1 ICBM & 9\% (alto)  & 20\% \\
\bottomrule
\end{tabular}
\caption{Imágenes BrainWeb utilizadas. Resolución 1 mm isótropo, formato NIfTI.}
\end{table}

El nivel de ruido \texttt{pn} representa el porcentaje de ruido Rician relativo a la
intensidad máxima del tejido blanco. El ruido Rician en imágenes de magnitud MRI es
asintóticamente equivalente a ruido Gaussiano para SNR $> 3$.

% ============================================================
\section{Resultados}
% ============================================================

\subsection{Vista global — Todos los filtros por nivel de ruido (axial)}

La Figura~\ref{fig:summary_axial} proporciona una vista consolidada del desempeño
comparativo de los tres filtros en los tres niveles de ruido, usando el corte axial central
del volumen.

\begin{figure}[H]
  \centering
  \includegraphics[width=\textwidth]{summary_axial.png}
  \caption{Comparación global: todos los filtros $\times$ todos los niveles de ruido.
           Corte axial central. Columnas: original, Mediana (r=1), Mediana Adaptativo,
           Wiener Adaptativo. Filas: pn3 (3\%), pn5 (5\%), pn9 (9\%).}
  \label{fig:summary_axial}
\end{figure}

Se observan las siguientes tendencias generales:
\begin{itemize}
  \item El Filtro de Mediana produce la imagen más nítida en pn3 y pn5, pero introduce
        un efecto de desenfoque notable en pn9.
  \item El Filtro de Mediana Adaptativo preserva mejor los bordes a bajo ruido y mantiene
        mayor detalle visible a alto ruido que el filtro de mediana clásico.
  \item El Filtro de Wiener Adaptativo produce una textura suave y homogénea a bajo y
        medio ruido, pero pierde detalle fino. A alto ruido (pn9) el efecto de suavizado
        es excesivo para la ventana 5×5×5 utilizada.
\end{itemize}

\subsection{Nivel de ruido pn3 (3\%) — Comparación completa}

\begin{figure}[H]
  \centering
  \includegraphics[width=\textwidth]{t1_icbm_normal_1mm_pn3_rf20_comparison.png}
  \caption{pn3 (3\%) — Comparación de los cuatro resultados en los tres planos
           anatómicos. Columnas: Original, Mediana (r=1), Mediana Adaptativo, Wiener
           Adaptativo.}
  \label{fig:pn3_full}
\end{figure}

\textbf{Análisis pn3:} A bajo nivel de ruido la imagen original ya exhibe buena calidad
diagnóstica. El Filtro de Mediana introduce un suavizado perceptible en las giro-surcos
cerebrales visibles en el plano sagital. El Filtro de Mediana Adaptativo conserva en mayor
medida los bordes de sustancia gris/blanca. El Filtro de Wiener Adaptativo produce la
imagen visualmente más suave, con pérdida de textura fina en el plano coronal — un indicio
de sobresuavizado para este nivel de ruido con ventana 5×5×5.

\subsubsection{Detalle: Mediana Adaptativo en pn3}

\begin{figure}[H]
  \centering
  \includegraphics[width=0.82\textwidth]{t1_icbm_normal_1mm_pn3_rf20_single_comparison.png}
  \caption{pn3 — Original vs.\ Filtro de Mediana Adaptativo ($S_{\max}=7$, variante NumPy).}
  \label{fig:pn3_amf}
\end{figure}

\subsubsection{Detalle: Wiener Adaptativo en pn3}

\begin{figure}[H]
  \centering
  \includegraphics[width=0.82\textwidth]{t1_icbm_normal_1mm_pn3_rf20_wiener_comparison.png}
  \caption{pn3 — Original vs.\ Filtro de Wiener Adaptativo (ventana 5×5×5, $\nu^2$ auto-estimada).}
  \label{fig:pn3_awf}
\end{figure}

\subsection{Nivel de ruido pn5 (5\%) — Comparación completa}

\begin{figure}[H]
  \centering
  \includegraphics[width=\textwidth]{t1_icbm_normal_1mm_pn5_rf20_comparison.png}
  \caption{pn5 (5\%) — Comparación de los cuatro resultados en los tres planos anatómicos.}
  \label{fig:pn5_full}
\end{figure}

\textbf{Análisis pn5:} A nivel de ruido medio el contraste entre los métodos se vuelve más
evidente. El Filtro de Mediana suprime efectivamente los puntos de ruido aislados que
comienzan a ser visibles, manteniendo contraste global. El Filtro de Mediana Adaptativo
muestra un balance favorable: suprime el ruido sin difuminar la delimitación córtex/CSF. El
Filtro de Wiener Adaptativo exhibe claramente su comportamiento de borde en el plano
sagital: las regiones de alta varianza (interfaces tejido/CSF) son preservadas con mayor
contraste, mientras las regiones homogéneas (tejido blanco profundo) son suavizadas de
manera efectiva.

\subsubsection{Detalle: Mediana Adaptativo en pn5}

\begin{figure}[H]
  \centering
  \includegraphics[width=0.82\textwidth]{t1_icbm_normal_1mm_pn5_rf20_single_comparison.png}
  \caption{pn5 — Original vs.\ Filtro de Mediana Adaptativo.}
  \label{fig:pn5_amf}
\end{figure}

\subsubsection{Detalle: Wiener Adaptativo en pn5}

\begin{figure}[H]
  \centering
  \includegraphics[width=0.82\textwidth]{t1_icbm_normal_1mm_pn5_rf20_wiener_comparison.png}
  \caption{pn5 — Original vs.\ Filtro de Wiener Adaptativo.}
  \label{fig:pn5_awf}
\end{figure}

\subsection{Nivel de ruido pn9 (9\%) — Comparación completa}

\begin{figure}[H]
  \centering
  \includegraphics[width=\textwidth]{t1_icbm_normal_1mm_pn9_rf20_comparison.png}
  \caption{pn9 (9\%) — Comparación de los cuatro resultados en los tres planos anatómicos.}
  \label{fig:pn9_full}
\end{figure}

\textbf{Análisis pn9:} A alto nivel de ruido el desempeño de los tres filtros se diferencia
claramente. El Filtro de Mediana logra recuperar estructuras anatómicas macroscópicas
(lobos, ventrículos) pero introduce artefactos de desenfoque severos en los planos sagital y
coronal — visible en la pérdida de definición de los giro-surcos y en la uniformización de la
textura cortical. El Filtro de Mediana Adaptativo mantiene más detalle que el Mediana
clásico, aunque a esta densidad de ruido la ventana máxima $S_{\max}=7$ comienza a ser
insuficiente para garantizar una mediana confiable en toda la imagen. El Filtro de Wiener
Adaptativo produce la imagen más suave de las tres pero con pérdida de detalle más
pronunciada: el suavizado global es excesivo y la varianza estimada $\nu^2$ se ve elevada
por el alto nivel de ruido, lo que deprime la ganancia en regiones de contraste moderado.

\subsubsection{Detalle: Mediana Adaptativo en pn9}

\begin{figure}[H]
  \centering
  \includegraphics[width=0.82\textwidth]{t1_icbm_normal_1mm_pn9_rf20_single_comparison.png}
  \caption{pn9 — Original vs.\ Filtro de Mediana Adaptativo.}
  \label{fig:pn9_amf}
\end{figure}

\subsubsection{Detalle: Wiener Adaptativo en pn9}

\begin{figure}[H]
  \centering
  \includegraphics[width=0.82\textwidth]{t1_icbm_normal_1mm_pn9_rf20_wiener_comparison.png}
  \caption{pn9 — Original vs.\ Filtro de Wiener Adaptativo.}
  \label{fig:pn9_awf}
\end{figure}

% ============================================================
\section{Análisis Comparativo}
% ============================================================

\subsection{Comparación entre los tres filtros}

\begin{table}[H]
\centering
\renewcommand{\arraystretch}{1.5}
\begin{tabularx}{\textwidth}{lXXX}
\toprule
\textbf{Aspecto} &
\textbf{Mediana (MF)} &
\textbf{Med. Adaptativo (AMF)} &
\textbf{Wiener Adaptativo (AWF)} \\
\midrule
Tipo de ventana &
Fija ($3\times3\times3$) &
Dinámica ($3\to5\to7$) &
Fija cúbica ($5\times5\times5$) \\

Preservación de bordes &
Moderada: suaviza bordes y ruido &
Alta: Nivel B conserva píxeles limpios &
Alta en bordes de alta varianza \\

Ruido de impulso &
Bueno a baja densidad; deteriora a alta &
Mejor tolerancia; $S_{\max}$ amplía la búsqueda &
Pobre: la mediana local queda sesgada \\

Ruido Gaussiano &
Bueno ($\hat{r}=1$ PSNR $\sim52$ dB a baja densidad) &
Aceptable en variante NumPy; la variante ITK introduce aclarado &
Óptimo para Gaussiano (estimador MMSE) \\

Sobresuavizado &
Moderado a radio pequeño &
Bajo (adapta la ventana) &
Presente cuando $\nu^2$ es sobreestimada \\

Velocidad (volumen T1) &
Muy rápido ($<$5 s con stride tricks) &
Rápido (var. ITK $\sim$30 s) / Muy lento (var. NumPy $>$10 min) &
Rápido ($<$10 s con \texttt{uniform\_filter}) \\

Implementación ITK &
\texttt{itk.MedianImageFilter} — nativo &
Aproximación: Median + GradientMagnitude + umbral &
Manual con NumPy/SciPy (no existe filtro nativo) \\

Parámetro clave &
\texttt{radius} &
$S_{\max}$ (\texttt{--max-window}) &
\texttt{window} (tamaño cúbico) y $\nu^2$ \\
\bottomrule
\end{tabularx}
\caption{Comparación de los tres filtros implementados.}
\label{tab:comparison}
\end{table}

\subsection{Desempeño por nivel de ruido}

\begin{table}[H]
\centering
\renewcommand{\arraystretch}{1.5}
\begin{tabular}{lccc}
\toprule
\textbf{Nivel} & \textbf{Mejor filtro} & \textbf{Balance borde/ruido} & \textbf{Observación clave} \\
\midrule
pn3 (3\%) & AWF / AMF & AMF superior en bordes &
  AWF sobresuaviza textura fina \\
pn5 (5\%) & AMF & AMF mejor balance global &
  Mediana introduce desenfoque leve \\
pn9 (9\%) & Mediana & Mediana recupera estructura macro &
  AWF pierde detalle; AMF limitado por $S_{\max}$ \\
\bottomrule
\end{tabular}
\caption{Mejor filtro por nivel de ruido (evaluación cualitativa).}
\end{table}

\subsection{Comparación con la implementación de Ali (2018)}

Ali (2018) evalúa los mismos tres filtros sobre imágenes MRI 2D con ruido Gaussiano y
sal y pimienta, reportando los siguientes valores PSNR representativos:

\begin{table}[H]
\centering
\renewcommand{\arraystretch}{1.3}
\begin{tabular}{lcccc}
\toprule
\textbf{Filtro} & \multicolumn{2}{c}{\textbf{Ruido Gaussiano (PSNR dB)}} &
                  \multicolumn{2}{c}{\textbf{Sal y pimienta (PSNR dB)}} \\
 & 10\% & 90\% & 10\% & 90\% \\
\midrule
Wiener Adaptativo & 43.21 & 36.53 & 45.25 & 33.47 \\
Mediana           & \textbf{51.98} & \textbf{45.54} & 61.82 & 39.13 \\
Mediana Adaptativo & 38.98 & 33.79 & \textbf{66.86} & \textbf{40.19} \\
\bottomrule
\end{tabular}
\caption{PSNR reportado por Ali (2018) para evaluación 2D. El Filtro de Mediana es
  superior para ruido Gaussiano; el Filtro de Mediana Adaptativo es superior para
  sal y pimienta.}
\end{table}

Nuestros resultados cualitativos en 3D son consistentes con estos hallazgos en cuanto a
la jerarquía de los filtros. Las diferencias observables se explican por:

\begin{enumerate}
  \item \textbf{Dimensionalidad}: nuestra implementación es 3D (cubos de vóxeles); Ali
        trabaja en 2D (matrices de píxeles). La vecindad 3D tiene acceso a información
        inter-corte que enriquece la estimación estadística.
  \item \textbf{Tipo de ruido BrainWeb}: el ruido de BrainWeb es Rician (no puramente
        Gaussiano ni puramente impulsivo). El AWF está teóricamente optimizado para
        Gaussiano; el AMF tiene ventaja conceptual sobre impulso.
  \item \textbf{Variante ITK del AMF}: la aproximación ITK del AMF (Median +
        GradientMagnitude + umbral) no es el algoritmo exacto de Ali. La variante NumPy
        sí lo es, pero es computacionalmente costosa para volúmenes 3D. Ali implementa
        el algoritmo exacto sobre 2D, lo que facilita tiempos razonables.
  \item \textbf{Estimación de $\nu^2$}: Ali no describe explícitamente cómo estima
        $\nu^2$ en el AWF. Nuestra implementación usa el promedio de varianzas locales
        (ecuación \ref{eq:nu2}), que puede sobreestimar $\nu^2$ cuando el volumen
        contiene regiones de alto contraste, deprimiendo la ganancia en esas zonas.
\end{enumerate}

% ============================================================
\section{Conclusiones}
% ============================================================

\begin{enumerate}

  \item \textbf{Filtro de Mediana (MF):} Método robusto y computacionalmente eficiente,
    superior para ruido Gaussiano de baja a media intensidad. La implementación 3D
    vectorizada con \texttt{sliding\_window\_view} es práctica para volúmenes clínicos
    estándar. Su principal limitación es el suavizado indiscriminado de bordes, que se
    acentúa a ruido alto (pn9) o con radios de ventana mayores.

  \item \textbf{Filtro de Mediana Adaptativo (AMF):} El mecanismo de ventana dinámica y
    el criterio de decisión por vóxel (Niveles A y B) le confieren mayor preservación de
    bordes y estructuras anatómicas finas respecto al filtro clásico. Es el método más
    versátil de los tres para ruido mixto de densidad moderada. La variante ITK ofrece
    velocidad a cambio de fidelidad al algoritmo exacto; para estudios que requieran el
    algoritmo estricto de Ali (2018) debe usarse la variante NumPy.

  \item \textbf{Filtro de Wiener Adaptativo (AWF):} Teóricamente óptimo bajo el criterio
    MMSE para ruido Gaussiano. En la práctica, su desempeño sobre BrainWeb depende
    críticamente de la estimación de $\nu^2$ y del tamaño de ventana. Para ruido bajo
    (pn3, pn5) ofrece una reducción de ruido suave y homogénea que puede ser deseable
    en aplicaciones de segmentación. Para ruido alto (pn9) el sobresuavizado compromete
    el detalle diagnóstico. Una calibración manual de $\nu^2$ o una ventana más pequeña
    mejorarían su desempeño en este escenario.

  \item \textbf{Elección del filtro:} En concordancia con Ali (2018), la elección debe
    guiarse por el tipo y densidad de ruido presente. Para ruido predominantemente
    Gaussiano (como en MRI clínica estándar), el Filtro de Mediana y el AWF bien
    calibrado son opciones sólidas. Para ruido de impulso o artefactos puntuales, el
    AMF con $S_{\max}$ adecuado es preferible.

  \item \textbf{Rol de ITK:} En los tres scripts, ITK cumple el rol de capa de
    entrada/salida (I/O) para los formatos médicos estándar (NIfTI, DICOM), manteniendo
    la metadata física (espaciado, origen, dirección) del volumen. La lógica de filtrado
    se implementa en NumPy/SciPy por mayor flexibilidad y control, dado que ITK no
    provee un filtro de Wiener adaptativo espacial nativo ni una implementación exacta
    del AMF de Ali (2018).

\end{enumerate}

% ============================================================
\section*{Referencias}
\addcontentsline{toc}{section}{Referencias}
% ============================================================

\begin{thebibliography}{9}

\bibitem{ali2018}
  Ali, H.\,M. (2018).
  \textit{MRI Medical Image Denoising by Fundamental Filters}.
  En: \textit{High-Resolution Neuroimaging — Basic Physical Principles and
  Clinical Applications}. IntechOpen.
  \url{https://doi.org/10.5772/intechopen.72427}

\bibitem{lin2013}
  Lin, L., Meng, X., \& Liang, X. (2013).
  Reduction of impulse noise in MRI images using block-based adaptive median filter.
  \textit{IEEE International Conference on Medical Imaging Physics and Engineering
  (ICMIPE)}, pp.\,132--134.

\bibitem{brainweb}
  Collins, D.\,L., et al. (1998).
  Design and construction of a realistic digital brain phantom.
  \textit{IEEE Transactions on Medical Imaging}, 17(3), 463--468.
  \url{https://doi.org/10.1109/42.712135}

\bibitem{itk}
  ITK Development Team. (2024).
  \textit{Insight Toolkit (ITK) v5.4}.
  \url{https://itk.org}

\bibitem{scipy}
  Virtanen, P., et al. (2020).
  SciPy 1.0: Fundamental algorithms for scientific computing in Python.
  \textit{Nature Methods}, 17, 261--272.
  \url{https://doi.org/10.1038/s41592-019-0686-2}

\end{thebibliography}

\end{document}